\section{Conclusiones}

A partir de los experimentos realizados se identificaron los siguientes hallazgos
principales:

\begin{enumerate}
    \item \textbf{Rendimiento:} en el entorno evaluado, los candidatos de retículos autentican entre 80 y 215 veces más lento que ECDSA, entre 5.85~ms (Ring-LWE) y 15.68~ms (LWR), en parte porque se comparan prototipos en Python con una implementación compilada. ML-ZKAuth, con 9.72~ms, queda entre los candidatos sobre anillos y los de matriz.

    \item \textbf{Memoria:} la estructura algebraica es el factor decisivo. Los candidatos de matriz asignan entre 16.4 y 18.7~MB por autenticación y requieren parámetros públicos de 6.6 a 7.6~MB, mientras que Ring-LWE y ML-ZKAuth asignan menos de 1~MB y sus claves ocupan pocos kilobytes.

    \item \textbf{Comunicación:} ML-ZKAuth transmite 4\,910 bytes por autenticación, entre 2.3 y 2.6 veces menos que los demás candidatos de retículos, aunque 18.4 veces más que ECDSA.

    \item \todoresult{hallazgo principal sobre seguridad}.

    \item \todoresult{hallazgo principal sobre despliegue}.

    \item \todoresult{hallazgo principal sobre el protocolo propuesto}.
\end{enumerate}

En conjunto, los resultados permiten caracterizar las diferentes alternativas no
mediante una clasificación única, sino a partir de los compromisos que cada una
establece entre seguridad, costo computacional, memoria, comunicación y complejidad
operacional.

Esta caracterización proporciona la base experimental para determinar en qué contextos
resulta apropiado utilizar cada familia de protocolos y permite evaluar de forma
homogénea la propuesta desarrollada en este trabajo.